\documentclass[11pt,a4paper]{article}
\usepackage[utf8]{inputenc}
\usepackage[french]{babel}
\usepackage[T1]{fontenc}
\usepackage{amsmath,amsfonts,amssymb}
\usepackage{geometry}
\geometry{margin=1.5cm}
\usepackage{tcolorbox}
\usepackage{multicol}

\title{\textbf{Fiche de Révision : Le Calcul Littéral (Bases)}}
\author{Mathématiques - Niveau 3\textsuperscript{e}}
\maDate{}

\begin{document}

\maketitle

    \begin{tcolorbox}[colback=blue!5!white,colframe=blue!75!black,title=Rappels Importants]
    \begin{itemize}
        \item \textbf{Réduire} : On ne peut additionner que les termes de la même "famille" (les $x^2$ ensemble, les $x$ ensemble, et les nombres seuls ensemble).
        \item \textbf{Distributivité simple} : $k \times (a + b) = k \times a + k \times b$
        \item \textbf{Double distributivité} : $(a + b)(c + d) = a \times c + a \times d + b \times c + b \times d$
        \item \textbf{Règle des signes} : $(+) \times (+) = (+)$ ; $(-) \times (-) = (+)$ ; $(+) \times (-) = (-)$
    \end{itemize}
    \end{tcolorbox}
    
    \section*{EXERCICE 1 : Réduction d'expressions}
    \textit{Simplifier les expressions suivantes au maximum.}
    
    \begin{multicols}{2}
    \begin{enumerate}
        \item $A = 3x + 5x$
        \item $B = 7x - 2x$
        \item $C = 4x + 3 + 2x - 1$
        \item $D = 5x^2 + 2x^2$
        \item $E = 8x^2 + 3x - 5x^2 + 4x$
        \item $F = 10 - 2x + 4 - 5x$
    \end{enumerate}
    \end{multicols}
    
    \section*{EXERCICE 2 : Multiplication et puissances}
    \textit{Attention à ne pas confondre addition et multiplication !}
    
    \begin{multicols}{2}
    \begin{enumerate}
        \item $G = x \times x$
        \item $H = 3x \times 2$
        \item $I = 4x \times 5x$
        \item $J = 2x \times (-3x)$
    \end{enumerate}
    \end{multicols}
    
    \section*{EXERCICE 3 : Calcul de valeur}
    \textit{On donne l'expression $E = 3x^2 - 5x + 2$. Calculer la valeur de $E$ pour :}
    
    \begin{enumerate}
        \item $x = 2$ \dotfill
        \item $x = 10$ \dotfill
        \item $x = 0$ \dotfill
    \end{enumerate}
    
    \newpage % Passage au verso
    
    \section*{EXERCICE 4 : Distributivité simple}
    \textit{Développer puis réduire les expressions suivantes.}
    
    \begin{enumerate}
        \item $K = 3(x + 4) = \dots\dots\dots\dots\dots\dots\dots\dots\dots\dots\dots\dots\dots\dots$
        \item $L = 5(2x - 3) = \dots\dots\dots\dots\dots\dots\dots\dots\dots\dots\dots\dots\dots\dots$
        \item $M = x(x + 7) = \dots\dots\dots\dots\dots\dots\dots\dots\dots\dots\dots\dots\dots\dots$
        \item $N = -2(4x - 5) = \dots\dots\dots\dots\dots\dots\dots\dots\dots\dots\dots\dots\dots\dots$
    \end{enumerate}
    
    \section*{EXERCICE 5 : Suppression de parenthèses}
    \textit{Rappel : Un signe "$-$" devant une parenthèse change tous les signes à l'intérieur.}
    
    \begin{enumerate}
        \item $O = 5 + (2x + 3) = \dots\dots\dots\dots\dots\dots\dots\dots\dots\dots\dots\dots\dots\dots$
        \item $P = 10 - (3x - 4) = \dots\dots\dots\dots\dots\dots\dots\dots\dots\dots\dots\dots\dots\dots$
    \end{enumerate}
    
    \section*{EXERCICE 6 : Double distributivité (Niveau +)}
    \textit{Développer et réduire chaque expression.}
    
    \begin{enumerate}
        \item $Q = (x + 2)(x + 3)$ \\
        $Q = x \times \dots + x \times \dots + 2 \times \dots + 2 \times \dots$ \\
        $Q = \dots\dots\dots\dots\dots\dots\dots\dots\dots\dots\dots\dots\dots\dots\dots\dots$
        
        \item $R = (x + 5)(2x - 1)$ \\
        $R = \dots\dots\dots\dots\dots\dots\dots\dots\dots\dots\dots\dots\dots\dots\dots\dots$
        
        \item $S = (2x - 3)(x - 4)$ \\
        $S = \dots\dots\dots\dots\dots\dots\dots\dots\dots\dots\dots\dots\dots\dots\dots\dots$
    \end{enumerate}
    
    \section*{EXERCICE 7 : Problème concret}
    \textit{On considère un rectangle de longueur $L = x + 5$ et de largeur $l = 3$.}
    
    \begin{enumerate}
        \item Exprimer le périmètre du rectangle en fonction de $x$. \\
        \textit{(Rappel : $P = 2 \times (L + l)$)} \\
        \dotfill
        \item Exprimer l'aire du rectangle en fonction de $x$ sous forme développée. \\
        \textit{(Rappel : $A = L \times l$)} \\
        \dotfill
    \end{enumerate}

\end{document}